\begin{figure}[t]
\small
\newjudge{Well-formed expressions}{\jtype{e}{\expty{\epsilon}{\alpha}}}
\vspace*{-.5em}
\infrule[T-Variable]
  { \Gamma = \dontcares, z\COLON\alpha, \rep{z'\COLON\dontcare}  \andalso  \rep{\same{z} \neq z'} }
  { \jtype{z}{\expty{z}{\alpha}} }

\infrule[T-Closure]
  { \mathit{distinct}\tupleb{x}  \andalso
    \rep{\jwftype[\same{\Delta}]{\tau}}  \andalso
    \jwftype{\omega}  \andalso
    \jtype[\Delta;\Gamma,\bar{x\COLON\tau}]{e}{\expty{\epsilon}{\rho}}  \andalso
    \jsubtype{\rho}{\omega} }
  { \jtype{((\bar{x\COLON\tau})\COLON\omega \Rightarrow e)}
             {\expty{((\bar{x\COLON\tau})\COLON\omega \Rightarrow \epsilon)}{\arrowtype{(\bar{\tau})}{\omega}}} }

\infrule[T-Apply]
  { \jtype{e}{\expty{\epsilon}{\arrowtype{(\bar{\alpha})}{\rho}}}  \andalso
    \rep{\jtype[\same{\Delta;\Gamma}]{e'}{\expty{\epsilon'}{\chi}}} \andalso
    \rep{\jsubtype[\same{\Delta}]{\chi}{\alpha}}  }
  { \jtype{e\apply(\bar{e'})}{\expty{\epsilon\apply(\bar{\epsilon'})}{\rho}} }

\infrule[T-Object]
  { \objectdecl{O}{\rep{\bindup{P}{\dontcares}}}{\rep{x\COLON\omega}}{\dontcares}{\dontcares} \in \Delta \\[3pt]
    \jwftype{O\obb{\tau}}  \andalso
    \rep{\jtype[\same{\Delta;\Gamma}]{e}{\expty{\epsilon}{\eta}}}  \andalso
    \rep{\jsubtype[\same{\Delta}]{\eta}{\underline{\Bigsubstb{\tau}{P}}\omega}} }
  { \jtype{O\obb{\tau}(\bar{e})}{\expty{O\obb{\tau}(\bar{\epsilon})}{O\obb{\tau}}} }

\infrule[T-Method]
  { \jtype{e}{\expty{\epsilon}{C}}  \andalso
    \rep{\jtype[\same{\Delta;\Gamma}]{e'}{\expty{\epsilon'}{\eta}}} \andalso
    \msav{\tuple{C,m}}{\tuple{C,\rep\eta}} = \methoddef{m}{}{\dontcare}{\rho}{\dontcare}
  }
  { \jtype{e.m(\bar{e'})}{\expty{\epsilon.m(\bar{\epsilon'})\COLON \rho}{\rho}} }

\newjudge{Most specific applicable visible method}{\msav{(C, m)}{\alpha} = D}
\vspace*{-.5em}

\infrule[MSAV]
  { \mostspecific{\applicable{\visible{(C, m)}}{\alpha}} = \setb{\tuple{d,\setb{\sigma}}} \andalso \LRcountofb{\tuple{d,\setb{\sigma}}}=1 \\
 \andalso \tuple{d',\setb{\sigma'}}\in\setb{\tuple{d,\setb{\sigma}}} \andalso
    \rep{\instantiate{\same{d'}}{\sigma'} = D'} \andalso D \in \setb{D'} }
  { \msav{(C, m)}{\alpha} = D }

\flushleft\textbf{Miscellaneous judgments:}\par
\flushright
\infax{
\fbox{$
\begin{array}{@{}c@{}}
\quad
\mostspecific{\setb{\tuple{d,\setb{\sigma}}}} = \setb{\tuple{d,\setb{\sigma}}} \andalso
\morespecific{d}{d'} \quad\null \\[2pt] \quad
\bigapplicable{\setb{d}}{\alpha} = \setb{\tuple{d,\setb{\sigma}}} \andalso
\visible{(C,m)} = \setb{d}  \quad\null \instantiate{d}{\sigma} = D
\end{array}
$}
}
\infax[Most-Specific]
  { 
\begin{array}{@{}l@{}}
\mostspecific{\setb{(d,\setb{\sigma})}} = \\
\quad\quad \Bigsetcomp{ (d',\setb{\sigma'}) }{ (d',\setb{\sigma'}) \in \setb{(d,\setb{\sigma})}, {\displaystyle\bigwedge_{(d'',\dontcare) \in \setb{(d,\setb{\sigma})}}\;\;\morespecific{d'}{d''}} }
\end{array}
 }

\infax[More-Specific]
  { \morespecific{d}{d'} = \jmorespecific{d}{d'} }

\infax[Applicable]
  { 
\begin{array}{@{}l@{}}
\bigapplicable{\setb{d}}{\alpha} =\\
\andalso
\LRsetcomp{ (d',\setb{\sigma}) }{ d' \in \setb{d},
\setb{\sigma} = \LRsetcomp{ \sigma' }{
\jleusing{\existstype{\verythinmathspace}{\alpha}}{\dom(d')}{\sigma'} }
\neq\emptyset }
\end{array}
 }

\infax[Visible]
  { \visiblemeth{C}{m} = \LRsetcomp{d}{d\in\allvisible{C}, \name{d} = m} }

\infax[Instantiate-Method]
  { \biginstantiate{\tuple{C, \methoddef{m}{\dontcares}{\rep{x\COLON\alpha}}{\rho}{e}}}{\sigma} = \bigtuple{C, \methoddef{m}{}{\rep{x\COLON\sigma\alpha}}{\sigma\rho}{\sigma e}} }


\vspace*{-1.5em}
\caption{Well-formed expressions}
\label{fig:wellformedexpr}
\end{figure}

\section{Overloading Resolution}\label{sec:resolve}
Method overloading is resolved both at compile time for type checking
and at run time for evaluation, and these
\textit{static and dynamic overloading resolutions} should be consistent.
For a method invocation expression, static overloading resolution
chooses an instance of the most specific method declaration that is
applicable to the static types of the method arguments;
dynamic overloading resolution then must select an instance of the
most specific method declaration that is applicable to the dynamic types
of the arguments and whose return type is a subtype of the static type of
the expression.
% In this section, we describe the static and dynamic
% overloading resolutions in detail, prove that they make type-sound
% selections, and discuss the design decisions for FGFV.

\subsection{Static Overloading Resolution}\label{sec:static}
As discussed in Section~\ref{sec:motiv:fgfv}, in order to make
type-sound choices, the static types of
method invocation expressions must be available at run time.
Thus, checking well-formedness of an expression $e$ produces its
type-annotated version $\epsilon$ and its type $\alpha$
as shown in Figure~\ref{fig:wellformedexpr}.  A \textit{type-annotated
internal expression} is just like an expression except that all the
method invocation expressions in it are annotated with their static
types:
\begin{center}
$\descsm{Type-annotated internal expression} \quad
\epsilon~::=~z~\mid~((\bar{x\COLON \alpha})\COLON \alpha \Rightarrow \epsilon) ~\mid~\epsilon\apply(\bar{\epsilon})~\mid~O\obb{\alpha}(\bar{\epsilon})~\mid~\epsilon.m(\bar{\epsilon})\COLON \alpha$
\end{center}

We focus on rule \rncite{T-Method} for method invocation, since
the rules for the other kinds of expressions are conventional.
For a method invocation $e.m(\bar{e'})$,
after checking well-formedness of expressions $e$ and $\bar{e'}$,
the rule performs static overloading resolution via rule \rncite{MSAV}
to find an instance of the most specific applicable visible method declaration,
and annotates the method invocation with the return type of
the chosen instance.  Rule \rncite{MSAV} first selects the most
specific applicable visible method declaration $d'$ with a set of
substitutions $\setb{\sigma'}$ using the static types of the method
arguments $\alpha$.
The unique existence of the most specific applicable visible 
method declaration is guaranteed by Theorem~\ref{thm:mostspecific}.
Then, the rule selects one instance $D$ out of instantiations of $d'$
with $\setb{\sigma'}$.

\subsection{Dynamic Overloading Resolution}\label{sec:dynamic}
When evaluating a method invocation expression, dynamic overloading
resolution selects a method declaration to call by using the dynamic
types of the arguments and the static type of the expression.
We describe the dynamic overloading resolution in a top-down manner
from the dynamic semantics rule for method invocation
% (Section~\ref{sec:dynamic:sem}) to the details of its helper judgments
% (Section~\ref{sec:dispatch} to Section~\ref{sec:solve}), present
(Section~\ref{sec:dynamic:sem}) to the details of its helper judgments
(Sections~\ref{sec:match} and~\ref{sec:solve}), present
a concrete example (Section~\ref{sec:ex}), and prove its correctness
(Section~\ref{sec:prop}).

\begin{figure}[t]
\leavevmode
\small
\begin{array}[t]{@{}lrl@{\hskip1ex}l@{}}
\desc{Class-declaration environment} & \Psi & ::= & \setb{\psi}\\
\desc{Ground type} & g & ::= & T\obb{g}~\mid~O\obb{g}~\mid~\tupleb{g}~\mid~\arrowtype{g}{g}~\mid~\Any~\mid~\uniontype{g}{g}~\mid~\intersectiontype{g}{g}\\
\desc{Value} & v & ::= & O\obb{g}(\bar{v})~\mid~((\bar{x \COLON g}) \COLON g \Rightarrow \epsilon)\\
\desc{Ilk} & k & ::= & O\obb{g}~\mid~\tupleb{g}~\mid~\arrowtype{g}{g}
\\
\desc{Evaluation context} & E & ::= & \emptycontext~\mid E\apply(\bar{\epsilon}) ~\mid v\apply(\bar{v},E,\bar{\epsilon}) ~\mid O\obb{g}(\bar{v},E,\bar{\epsilon}) ~\mid E.m(\bar{\epsilon})\COLON g ~\mid v.m(\bar{v},E,\bar{\epsilon})\COLON g\\
\desc{Redex} & R & ::= & v\apply(\bar{v}) ~\mid v.m(\bar{v})\COLON g
\end{array}

\vspace*{.5em}
\newjudge{Dynamic semantics}{\jevalstep{E\context{R}}{E\context{\epsilon}}}
\vspace*{-.5em}

\infrule[R-Method]
  {   v'' = O\obb{g}(\bar{v})
\andalso
\objectdecl{O}{\dontcares}{\rep{x\COLON\dontcare}}{\dontcares}{\dontcares}\in\denv
\andalso
\listofb{d} = \topologicalsort[\denv]{\visible[\denv]{\tuple{O\obb{g},m}}}
\\
(\dontcare, \methoddef{m}{}{\rep{x'\COLON \dontcare}}{\dontcare}{\epsilon}) =
 \firstdispatchable[\denv]{\listofb{d}}
{\tuple{\ilk{v''}, \bar{\ilk{v'}}}}
{g'}
}
  { \jevalstep{E\context{v''.m(\bar{v'})\COLON g'}}{E\context{ \bigl[v''/\kwd{self},\bar{v/x},\bar{v'/x'}\bigr]\epsilon}} }

%\vspace*{-1.5em}
\vspace*{-1em}
\newjudge{First dispatchable method instance}{\firstdispatchable[\denv]{\listofb{d}}{k}{g} = D}
\vspace*{-.5em}

\infrule[First-Dispatchable]
  { \listofb{d} = \listof{\bar{d'},d'',\dontcares}  \andalso
\rep{\neg({ \jdispatch[\same{\denv}]{\same{k}}{\same{g}}{d'}{\dontcare} })}
\\[2pt]
\jdispatch{k}{g}{d''}{\sigma}
\andalso
D = \instantiate{d''}{\sigma}
 }
  { \firstdispatchable[\denv]{\listofb{d}}{k}{g} = D }

\newjudge{Sorting of method declarations}{ \topologicalsort{\setb{d}} = \listofb{d} }
\vspace*{-.5em}

\infrule[Topological-Sort]
  { \setb{d}=\setb{d'} \andalso
    \rep{ \neg(\morespecific{d'_j}{d'_i}) }[1 \leq i < j \leq \countofb{d'}] }
  { \topologicalsort{\setb{d}} = \listofb{d'} }

%\vspace*{-.5em}
\newjudge{Dispatch semipredicate}{ \jdispatch{k}{g}{d}{\sigma} }
\vspace*{-.5em}

\infrule[Dispatch]
  {
 \arrow(d) = {\foralltypeb{K}{\arrowtype{\alpha}{\rho}}}
    \andalso
    \jmatch{k}{\obb{K}}{\alpha}{\bindco}{\biglistofb{\listofb{l}}}{\biglistofb{\listofb{u}}}
    \\[1pt]
\andalso\andalso\andalso\andalso\andalso\andalso\quad
    \jmatch{g}{\obb{K}}{\rho}{\bindcontra}{\biglistofb{\listofb{l'}}}{\biglistofb{\listofb{u'}}}
    \\[1pt]
    \jsolve
	  {\obb{K}}
	  {\cartesiantypeunion{\rep{K}}\Biglistof{\biglistofb{\listofb{l}},\biglistofb{\listofb{l'}}}}
	  {\cartesiantypeintersection{\rep{K}}\Biglistof{\biglistofb{\listofb{u}},\biglistofb{\listofb{u'}}}}
	  {\sigma} }
  { \jdispatch{k}{g}{d}{\sigma} }

\vspace*{-1.5em}
\caption{Dynamic semantics of method invocation expressions}
\label{fig:dynamic}
\end{figure}


\subsubsection{Dynamic Semantics}\label{sec:dynamic:sem}~\par
Dynamic overloading resolution uses additional metavariables as
in Figure~\ref{fig:dynamic}.  A class-declaration environment $\Psi$ is
a set of class declarations.  A ground type $g$ is a dynamically
available type, which is an internal type that does not contain type
parameters nor $\Bottom$; our type system requires that both static and
dynamic type inference algorithms instantiate type parameters with
non-$\Bottom$ types, and that all the type parameters are instantiated with
ground types at run time.  We call the run-time type of
a value the \textit{ilk} of the value; the ilk of an object instance
$O\obb{g}\tupleb{v}$ is $O\obb{g}$, and the ilk of a function
$((\bar{x \COLON g}) \COLON g' \Rightarrow \epsilon)$ is
$\tuple{\barearrowtype{(\bar{g})}{g'}}$.  Similarly to domain types,
we write $\tupleb{g}$ to denote ilks of multiple function parameters.
Evaluation contexts $E$ and redexes $R$ are conventional.

For a method invocation, rule \rncite{R-Method} seeks the most
specific method declaration with a \textit{dispatchable instance} that
is applicable to the ilk of the arguments and whose return type is a
subtype of the annotated static return type.
We say a method declaration is \textit{dispatchable} to the ilk of
arguments and the static return type if it has a dispatchable instance.
For the sake of efficiency, rule \rncite{R-Method} first topologically sorts
all the visible method declarations in a most-to-least specific order.
% Since the specificity of method declarations is statically determinable,
% we can also topologically sort the method declarations at compile time
% to make dynamic overloading resolution more efficient.
(This order is statically determinable, so the sorting can be done at compile time.)
% When the sorted list $\listofb{d}$ is given, 
% the rule then checks the \textit{dispatch semipredicate} 
% for each method declaration from the beginning of the list 
The rule then checks the \textit{dispatch semipredicate} 
for each method declaration in the sorted list 
until it finds the first method declaration $d''$ dispatchable to the given
ilk of the arguments $k$ and the static return type $g$ with a dispatchable
instance $D=\instantiate{d''}{\sigma}$.


%\subsubsection{Dispatch Semipredicate}\label{sec:dispatch}~\par
The dispatch semipredicate defined in
% Figure~\ref{fig:dispatch}
rule \rncite{Dispatch}
checks if
a given method declaration $d$ is dispatchable to a given ilk $k$ and
a static return type $g$.  The dispatch semipredicate is satisfied using
substitution $\sigma$ if $d$ has a dispatchable instance (obtained by deleting the method type parameters of $d$
and applying $\sigma$ to its argument types, return type, and body).
Metavariable $\lbound$ and $a$ range over lower type bounds, and
$\ubound$ and $b$ range over upper type bounds.
We write $\listofb{\lbound}$ to denote ordered lists.
The dispatch semipredicate consists of three steps: (1)~computing type bounds
$\biglistofb{\listofb{\lbound}}$ and $\biglistofb{\listofb{\ubound}}$
that will make the domain type of 
the method declaration $d$ be a supertype of the ilk $k$, so that $d$ will be applicable to $k$,
(2)~computing type bounds
$\biglistofb{\listofb{\lbound'}}$ and $\biglistofb{\listofb{\ubound'}}$
that make the static return type $g$ a supertype of the return type of $d$
to make sure that the dynamic dispatch preserves type,
and (3)~solving those type bounds to get a valid substitution $\sigma$.
(If any of these steps fails, the semipredicate is false: the declaration is not applicable.)
The computed type bounds
% for the method type parameters
are either
ground types or $\Bottom$.  The dispatch semipredicate computes type
bounds by type matching as described in the next section.


\subsubsection{Computing Bounds by Type Matching}\label{sec:match}~\par
\begin{figure}[t]
\small
\vspace*{-1em}
\newjudge{Type matching a ground type to an internal type}{\jmatch{g}{\obb{K}}{\alpha}{V}{\biglistofb{\listofb{\lbound}}}{\biglistofb{\listofb{\ubound}}}}
\vspace*{-.5em}

\infrule[M-Param-Co]
  { \hbox{$g$ is not an intersection type} \andalso \hbox{$g$ is not a union type} \\[1pt]
    \rep{\lbound = \ifexpr{(P = \same{P_i})}{\same{g}}{\Bottom}}
\andalso
    \rep{\ubound = \Any} \andalso
    1 \leq i \leq \countofb{P} = \countofb{\ubound}
 }
  { \jmatch{g}{\obb{\lobindup{\rep{\chi}}{P}{\rep{\eta}}}}{P_i}{\bindco}{\biglistof{\listofb{\lbound}}}{\biglistof{\listofb{\ubound}}} }

\infrule[M-Tuple-Tuple-Success]
  { \rep{\jmatch[\same{\denv}]{g}{\same{\obb{K}}}{\alpha}{\same{V}}{\biglistofb{\listofb{\lbound}}}{\biglistofb{\listofb{\ubound}}} } }
  { \jmatch{\tupleb{g}}{\obb{K}}{\tupleb{\alpha}}{V}{\cartesiantypeunion{\rep{K}} \Biglistofb{\biglistofb{\listofb{\lbound}}}}{\cartesiantypeintersection{\rep{K}} \Biglistofb{\biglistofb{\listofb{\ubound}}}} }

\vspace*{-1.7em}
\infrule[M-Tuple-Tuple-Failure]
  { \countofb{g} \neq \countofb{\alpha} }
  { \jmatch{\tupleb{g}}{\dontcare}{\tupleb{\alpha}}{\dontcare}{\biglistof{}}{\biglistof{}} }

% \infrule[M-Arrow-Arrow]
%   { \jmatch{g}{\obb{K}}{\alpha}{\flip{V}}{\biglistofb{\listofb{\lbound}}}{\biglistofb{\listofb{\ubound}}} \\
%     \jmatch{h}{\obb{K}}{\rho}{V}{\biglistofb{\listofb{\lbound'}}}{\biglistofb{\listofb{\ubound'}}} }
%   { \jmatch{\arrowtype{g}{h}}{\obb{K}}{\arrowtype{\alpha}{\rho}}{V}{\cartesiantypeunion{\rep{K}} \Biglistof{\biglistofb{\listofb{\lbound}},\biglistofb{\listofb{\lbound'}}}}{\cartesiantypeintersection{\rep{K}} \Biglistof{\biglistofb{\listofb{\ubound}},\biglistofb{\listofb{\ubound'}}}} }

% \vspace*{-2em}
% \infrule[M-Union-Whatever-Co]
%   { \alpha \text{ is not an intersection type} \\
% \rep{\jmatch[\same{\denv}]{g}{\same{\obb{K}}}{\same{\alpha}}{\bindco}{\biglistofb{\listofb{\lbound}}}{\biglistofb{\listofb{\ubound}}}} }
%   { \jmatch{\naryuniontypeb{g}}{\obb{K}}{\alpha}{\bindco}{\cartesiantypeunion{\rep{K}} \Biglistofb{\biglistofb{\listofb{\lbound}}}}{\cartesiantypeintersection{\rep{K}} \Biglistofb{\biglistofb{\listofb{\ubound}}}} }

\infrule[M-Union-Whatever-Contra]
  { \hbox{$\alpha$ is not an intersection type} \andalso \hbox{$\alpha$ is not a union type} \\[1pt]
    \rep{\jmatch[\same{\denv}]{g}{\same{\obb{K}}}{\same{\alpha}}{\bindcontra}{\biglistofb{\listofb{\lbound}}}{\biglistofb{\listofb{\ubound}}}} }
  { \jmatch{\naryuniontypeb{g}}{\obb{K}}{\alpha}{\bindcontra}{\Biglistofb{\rep{\listofb{\lbound}}}}{\Biglistofb{\rep{\listofb{\ubound}}}} }
\vspace*{-1.5em}
\caption{Type matching a ground type to an internal type (selected rules)}
\label{fig:match}
\end{figure}

The \textsf{match} rules in Figure~\ref{fig:match} compute lower and upper
type bounds $\biglistofb{\listofb{\lbound}}$ and
$\biglistofb{\listofb{\ubound}}$ for method type parameters in $\bar{K}$
by matching a ground type $g$ against an internal type $\alpha$ according
to a variance mark $V$.
When $V$ is respectively $\bindco$, $\bindcontra$, or $\bindinv$, the rules compute 
the type bounds that make $g$ respectively a subtype, a supertype, or an equivalent type of $\alpha$.
For example, as in~\rncite{M-Param-Co}, if a ground type $g$ is
matched against a method type parameter $P_i$ according to $\bindco$,
$g$ becomes the lower bound of $P_i$ so that $g$ is a subtype of $P_i$
as long as $P_i$ satisfies its type bounds $g \subtypeof P_i \subtypeof \Any$.

The \textsf{match} rules compute lower and upper type bounds
in the form of two lists of lists
$\biglistofb{\listofb{\lbound}}$ and $\biglistofb{\listofb{\ubound}}$
where $\countofb{\listofb{\lbound}} = \countofb{\listofb{\ubound}}$ and
$\rep{\countofb{K} = \countofb{\lbound_i} = \countofb{\ubound_i}}[i]$.
Matching may produce multiple solutions; each inner list is one possible solution.
Multiple solutions come from cases with union types and
intersection types as in rule~\rncite{M-Union-Whatever-Contra}, which
matches a union ground type $\naryuniontypeb{g}$ against $\alpha$
according to $\bindcontra$ and computes type bounds that make
$\naryuniontypeb{g}$ a supertype of $\alpha$.  Because a union type
$\naryuniontypeb{g}$ is a supertype of $\alpha$ if any component $g_i$
is a supertype of $\alpha$, multiple components may be supertypes of
$\alpha$, which may lead to multiple type bounds for each of those
components.  Since we cannot decide which one of the multiple type bounds
is the one eventually satisfying all the constraints, the \textsf{match}
rules compute all possible solutions, which are then winnowed in the solving step.
To collect all the bounds, the \textsf{match} rules cover failure
cases as well so that the rules do not stop collecting bounds even
when one component fails type matching.  For example,
rule~\rncite{M-Tuple-Tuple-Failure} specifies a failure case due to
size mismatch of tuple types; rather than actually failing,
it succeeds but delivers empty lists of solutions.
The entire \textsf{match} rules are defined in Appendix~\ref{app:match}.

\begin{figure}[t]
\small
\newjudge{Transposing reduction on list of lists}{
    \transposedtypeunion\biglistofb{\listofb{\alpha}} = \listofb{\alpha}
    \quad
    \transposedtypeintersection\biglistofb{\listofb{\alpha}} = \listofb{\alpha} }

\infax{ \transposedtypeunion \LRlistof{\rep{\LRlistof{\rep{\alpha_{i\,j}}[j]}}[i]}
        = {\LRlistof{\rep{\bigtypeunion \LRlistof{\rep{\alpha_{i\,j}}[i]}}[j]}}
        \andalso
        \transposedtypeintersection \LRlistof{\rep{\LRlistof{\rep{\alpha_{i\,j}}[j]}}[i]}
        = {\LRlistof{\rep{\bigtypeintersection \LRlistof{\rep{\alpha_{i\,j}}[i]}}[j]}}
}

\newjudge{Reduction of a cross product}{
    \cartesiantypeunion{\rep{K}}\Biglistofb{\biglistofb{\listofb{\alpha}}} = \biglistofb{\listofb{\alpha}}
    \quad
    \cartesiantypeintersection{\rep{K}}\Biglistofb{\biglistofb{\listofb{\alpha}}} = \biglistofb{\listofb{\alpha}} }

\infax{ \cartesiantypeunion{\rep{K}}\Biglistof{} = \biglistof{\listof{\rep{\Bottom}[\countofb{K}]}} }

\infax{ \cartesiantypeunion{\rep{K}}\Biglistof{\rep{\biglistofb{\listofb{\alpha}}},\biglistofb{\listofb{\eta}}} =
    \LRlistof{\rep{\rep{\transposedtypeunion \biglistof{\listof{\rep{\gamma_{i\,k}}[k]},
                                                        \listof{\rep{\eta_{j\,k}}[k]}}}[j]}[i]}\quad
    \hbox{$\mathbf{where}$}\quad
          \biglistofb{\listofb{\gamma}} = \cartesiantypeunion{\rep{K}} \Biglistofb{\biglistofb{\listofb{\alpha}}}
  }

\infax{ \cartesiantypeintersection{\rep{K}}\Biglistof{} = \biglistof{\listof{\rep{\Any}[\countofb{K}]}} }

\infax{ \cartesiantypeintersection{\rep{K}}\Biglistof{\rep{\biglistofb{\listofb{\alpha}}},\biglistofb{\listofb{\eta}}} =
    \LRlistof{\rep{\rep{\transposedtypeintersection \biglistof{\listof{\rep{\gamma_{i\,k}}[k]},
                                                               \listof{\rep{\eta_{j\,k}}[k]}}}[j]}[i]}\quad
    \hbox{$\mathbf{where}$}\quad
          \biglistofb{\listofb{\gamma}} = \cartesiantypeintersection{\rep{K}} \Biglistofb{\biglistofb{\listofb{\alpha}}}
  }

\caption{Cartesian product reduction operators}
\label{fig:redop}
\end{figure}

% Lists of solutions are combined (in the style of a Cartesian product)
% after \textsf{match} rules match component types of a given type inductively.
Lists of solutions for component types are combined in the manner of a Cartesian product
as defined in Figure~\ref{fig:redop}.
For instance,
rule~\rncite{M-Tuple-Tuple-Success} reduces
$\cartesiantypeunion{\rep{K}}\Biglistof{\rep{\biglistofb{\listofb{\lbound_i}}}[i]}$
\Big($\cartesiantypeintersection{\rep{K}}\Biglistof{\rep{\biglistofb{\listofb{\ubound_i}}}[i]}$\Big)
by computing componentwise unions (intersections) for all the combinations of
choosing one type bound from each $i$-th list of the type bounds.
Hence, inductively matching $n$ components whose $i$-th component has
$m_i$ type bounds produces a list of $m_1 \times \cdots \times m_n$
type bounds (thus if matching any component fails and produces $\biglistof{}$,
the whole type matching fails and results in $\biglistof{}$).
For example,\\[-.5em]
\[
\begin{array}{l}
\cartesiantypeunion{\rep{K}}\Biglistof{\biglistof{\listof{A,B}, \listof{C,D}}, \biglistof{\listof{E,F}, \listof{G,H}, \listof{I,J}}}
=\\
 \biglistof{\listof{\bareuniontype{A}{E},\bareuniontype{B}{F}},
\listof{\bareuniontype{A}{G},\bareuniontype{B}{H}},
\listof{\bareuniontype{A}{I},\bareuniontype{B}{J}},
\listof{\bareuniontype{C}{E},\bareuniontype{D}{F}},
\listof{\bareuniontype{C}{G},\bareuniontype{D}{H}},
\listof{\bareuniontype{C}{I},\bareuniontype{D}{J}}}.
\end{array}
\]


\subsubsection{Solving the Bounds}\label{sec:solve}~\par
\begin{figure}[t]
\small
\newjudge{Solving type bounds}{ \jsolve{\obb{K}}{\biglistofb{\listofb{\lbound}}}{\biglistofb{\listofb{\ubound}}}{\sigma} }
\vspace*{-1em}

\infax[Solve-Empty]
  { \jsolve{\ob{}}{\biglistof{\listof{}}}{\biglistof{\listof{}}}{\emptysubst} }

\vspace*{-1em}
\infrule[Solve-Step]
  { \Bigcountofb{\listof{\rep\lbound,\lbound'}}=\Bigcountofb{\listof{\rep\ubound,\ubound'}}=n\geq 1
    \andalso 1\leq i\leq n
    \andalso g'=\intersectiontype{\ubound'_i}{\bigtypeintersection\setb{\eta'}}
    \andalso \jsubtype[\denv]{\lbound'_i}{g'}  \\
    \rep{\jmatch[\same{\denv}]{\same{g'}}{\same{\obb{K}}}{\chi'}{\bindcontra}{\biglistofb{\listofb{\lbound''}}}{\biglistofb{\listofb{\ubound''}}}} \\
    \jsolve
	   {\obb{K}}
	   {\cartesiantypeunion{\rep{K}}\Biglistof{\rep{\biglistofb{\listofb{\lbound''}}},\biglistof{\listofb{\lbound_{\same{i}}}}}}
	   {\cartesiantypeintersection{\rep{K}}\Biglistof{\rep{\biglistofb{\listofb{\ubound''}}},\biglistof{\listofb{\ubound_{\same{i}}}}}}
	   {\substb{g}{P}} }
  { \jsolve
       {\ob{\rep{K},\lobindup{\rep{\chi'}}{P'}{\rep{\eta'}}}}
	   {\biglistofb{\listof{\rep{\lbound},\lbound'}}}
	   {\biglistofb{\listof{\rep{\ubound},\ubound'}}}
	   {\substbx{g}{P}{g'}{P'}} }
\vspace*{-1.5em}
\caption{Solving type bounds}
\label{fig:solve}
\end{figure}

The last step of the dispatch semipredicate is to solve the computed type
bounds for method type parameters as in Figure~\ref{fig:solve}.
The \textsf{solve} rules compute a valid substitution $\sigma$ for all the
method type parameters in $\bar{K}$, which satisfies the declared
bounds in ${\bar{K}}$ and the computed bounds from type matching
$\biglistofb{\listofb{\lbound}}$ and $\biglistofb{\listofb{\ubound}}$.
The solving step succeeds if any of the type bounds leads to a valid
substitution.  For declared type bounds for method type parameters
$\rep{P}, P'$ and their computed type bounds
$\biglistofb{\listof{\rep\lbound,\lbound'}}$ and
$\biglistofb{\listof{\rep\ubound,\ubound'}}$ from type matching,
rule~\rncite{Solve-Step} nondeterministically selects $i$-th type bounds
$\listof{\rep{\lbound_{\same{i}}}, \lbound'_i}$ and
$\listof{\rep{\ubound_{\same{i}}}, \ubound'_i}$ that successfully compute
a valid substitution.  Starting from the rightmost type parameter
$P'$, the rule computes a substitution $\subst{g'}{P'}$ by taking the
intersection of the upper bounds of $P'$.  Because $g'$ is a ground type,
the solving step never instantiates method type parameters with $\Bottom$.
The \textsf{match} premises propagate the choice of $g'$ for $P'$ to the
declared lower bounds of $P'$ to collect additional constraints on
$\rep{P}$.  Then, the rule combines the additional constraints and the
originally computed bounds, and inductively solves the combined bounds
to compute the substitution $\substb{g}{P}$.  From this, it produces a valid
substitution $\substbx{g}{P}{g'}{P'}$.

Note that the rule selects $i$-th type bounds nondeterministically while
it checks type parameters deterministically from the rightmost one.
Because the scope of a well-formed method type parameter extends to the
right and method type parameters can appear only in lower bounds,
for well-formed type parameters $\rep{\lobindup{\bar\chi}{P}{\bar\eta}}$,
the upper bounds $\bar{\setb{\eta}}$ do not contain any of $\bar{P}$, and
$P_j$ does not appear in the lower bounds of $P_i$ for $i\le j$.
Thus, the choice of $g_i$ for $P_i$ does not affect the upper
bounds of $P_j$, which allows us to compute a substitution in the
right-to-left order by taking intersections of upper bounds.
Since the choice of $g_j$ for $P_j$ may place additional constraints on
$P_i$'s that appear in the lower bounds $\setb{\chi_{\same{j}}}$ of $P_j$,
\rncite{Solve-Step} propagates the choice of $g_j$ for $P_j$ to the
lower bounds and computes additional constraints, to ensure
that $g_j$ is a supertype of all $\setb{\chi_{\same{j}}}$.


\subsubsection{Example}\label{sec:ex}~\par
Let us revisit the dynamic overloading resolution process with a concrete example.
Suppose we have two covariant traits $\kwd{List}$ and $\kwd{SortedList}$
such that $\kwd{SortedList}\ob{P} \extends \kwd{List}\ob{P}$,
two objects $\kwd{Cons}$ and $\kwd{SortedCons}$ such that
$\kwd{Cons}\ob{P} \extends \kwd{List}\ob{P}$ and
$\kwd{SortedCons}\ob{P} \extends \kwd{SortedList}\ob{P}$, and
a method $\mathit{append}$ that concatenates two lists.
Consider the following method declarations:
\[
\begin{array}{l}
\mathit{append}\ob{P, \bd{P} \extends Q}(x \COLON \kwd{SortedList}\ob{P}, y \COLON \kwd{SortedList}\ob{Q}) \COLON \kwd{SortedList}\ob{Q} = \dontcare\\
\mathit{append}\ob{P,\bd{P} \extends Q}(x \COLON \kwd{List}\ob{P}, y \COLON \kwd{List}\ob{Q}) \COLON \kwd{List}\ob{Q} = \dontcare
\end{array}
\]
and call them $d_1$ and $d_2$ in order.  This is a valid overloading
since it satisfies the overloading rules with $d_1$ being more specific
than $d_2$.  Suppose we have three monomorphic classes $C_1$, $C_2$,
and $C_3$ such that $C_1 \subtypeof C_3$ and $C_2  \subtypeof C_3$, and
consider the following method invocation expression:
\begin{center}
$\mathit{append}(l_1, l_2)$
\end{center}
with two lists $l_1$ and $l_2$ whose static types are both
$\kwd{List}\ob{C_3}$ and whose ilks are $\kwd{SortedCons}\ob{C_1}$ and
$\kwd{Cons}\ob{C_2}$, respectively.  Then, type checking this expression
would produce the internal expression $\mathit{append}(l_1, l_2) \COLON \kwd{List}\ob{C_3}$.

Now, let us see how this method call is dynamically dispatched.
Given the ilk of the argument tuple
$\tuple{\kwd{SortedCons}\ob{C_1},\kwd{Cons}\ob{C_2}}$,
the static return type $\kwd{List}\ob{C_3}$, and the topologically sorted
list of method declarations $\listof{d_1, d_2}$, dynamic overloading resolution
performs the dispatch semipredicate starting from the first declaration $d_1$
in the list.
The dispatch semipredicate on $d_1$ would succeed if it computed a valid substitution.
However, it fails because its first type matching step fails, producing an empty list of solutions:
\medskip

\begin{figure}[h]
\small
\vspace*{-2em}
\infrule
  {
    \jmatch{\kwd{SortedCons}\ob{C_1}}{\ob{\bar{K}}}{\kwd{SortedList}\ob{P}}{\bindco}{\biglistof{\listof{C_1, \Bottom}}}{\biglistof{\listof{\Any, \Any}}} \\[2pt]
    \jmatch{\kwd{Cons}\ob{C_2}}{\ob{\bar{K}}}{\kwd{SortedList}\ob{Q}}{\bindco}{\biglistof{}}{\biglistof{}}
  }
  { \jmatch{\tuple{\kwd{SortedCons}\ob{C_1}, \kwd{Cons}\ob{C_2}}}{\ob{\bar{K}}}{\tuple{\kwd{SortedList}\ob{P}, \kwd{SortedList}\ob{Q}}}{\bindco}{\biglistof{}}{\biglistof{}} }
\vspace*{-2em}
\end{figure}

\medskip
\noindent
where $\bar{K} = \bd{\Bottom} \extends P \extends \bd{\Any}, \bd{P} \extends Q \extends \bd{\Any}$.
While the first component matching succeeds, the second component matching
fails since $\kwd{Cons}\ob{C_2}$ is not a subtype of $\kwd{SortedList}\ob{Q}$.
Although matching the static return type against the return type
of $d_1$ succeeds, the combined type bounds for the solving step are the empty
lists, which leads to the failure of the dispatch semipredicate:

\medskip
\begin{figure}[h]
\small
\vspace*{-2em}
\infrule
  { \arrow(d_1) = {\foralltype{\bar{K}}{\arrowtype{\tuple{\kwd{SortedList}\ob{P}, \kwd{SortedList}\ob{Q}}}{\kwd{SortedList}\ob{Q}}}}
    \\[2pt]
    \jmatch{\tuple{\kwd{SortedCons}\ob{C_1}, \kwd{Cons}\ob{C_2}}}{\ob{\bar{K}}}{\tuple{\kwd{SortedList}\ob{P}, \kwd{SortedList}\ob{Q}}}{\bindco}{\biglistof{}}{\biglistof{}}
    \\[2pt]
\jmatch{\kwd{List}\ob{C_3}}{\ob{\bar{K}}}{\kwd{SortedList}\ob{Q}}{\bindcontra}{\biglistof{\listof{\Bottom, \Bottom}}}{\biglistof{\listof{\Any, C_3}}}
    \\[2pt]
\jsolve{\ob{\bar{K}}}{\biglistof{}}{\biglistof{}}{\failed}
  }
  { \jdispatch{\tuple{\kwd{SortedCons}\ob{C_1}, \kwd{Cons}\ob{C_2}}}{\kwd{List}\ob{C_3}}{d_1}{\failed} }
\vspace*{-2em}
\end{figure}

\medskip
Dynamic overloading resolution then moves on to the next method declaration
in the list, $d_2$, and performs the dispatch semipredicate on it.
In this case, both type matching steps succeed, and the overall result
of the dispatch semipredicate depends on the result of the solving step:

\begin{figure}[h]
\small
\vspace*{-1.5em}
\infrule
  { \arrow(d_2) = {\foralltype{\bar{K}}{\arrowtype{\tuple{\kwd{List}\ob{P}, \kwd{List}\ob{Q}}}{\kwd{List}\ob{Q}}}}
    \\%[2pt]
    \jmatch{\tuple{\kwd{SortedCons}\ob{C_1}, \kwd{Cons}\ob{C_2}}}{\ob{\bar{K}}}{\tuple{\kwd{List}\ob{P}, \kwd{List}\ob{Q}}}{\bindco}{\biglistof{\listof{C_1, C_2}}}{\biglistof{\listof{\Any, \Any}}}
    \\%[2pt]
\jmatch{\kwd{List}\ob{C_3}}{\ob{\bar{K}}}{\kwd{List}\ob{Q}}{\bindcontra}{\biglistof{\listof{\Bottom, \Bottom}}}{\biglistof{\listof{\Any, C_3}}}
    \\%[2pt]
\jsolve{\ob{\bar{K}}}{\biglistof{\listof{C_1, C_2}}}{\biglistof{\listof{\Any, C_3}}}{\sigma}
  }
  { \jdispatch{\tuple{\kwd{SortedCons}\ob{C_1}, \kwd{Cons}\ob{C_2}}}{\kwd{List}\ob{C_3}}{d_2}{\sigma} }
\vspace*{-2em}
\end{figure}

\noindent
We start solving the type bounds from the rightmost type parameter $Q$.
Since we have a single pair of type bounds, that is,
$\bigcountof{\listof{C_1, C_2}} = \bigcountof{\listof{\Any, C_3}} = 1$,
there is no need for a nondeterministic choice in rule \rncite{Solve-Step}.
The substitution for $Q$ is the intersection of the computed upper bound
and the declared upper bound: $\bareintersectiontype{C_3}{\Any} = C_3$.
This choice of substitution $C_3$ for $Q$ has to be propagated by matching it
against the declared lower bound of $Q$, which is $P$, to make sure that
$C_3$ is a supertype of $P$.
We then combine the type bounds obtained from propagation and the originally
given type bounds, and inductively solve the combined bounds to get the
substitution for $P$, which results in $\subst{C_3}{P}$ by again taking the
intersection of the upper bounds: $\bareintersectiontype{C_3}{\Any} = C_3$.
Thus, the solving step succeeds with a valid substitution
$\substxx{C_3}{P}{C_3}{Q}$:

\smallskip
\begin{figure}[h]
\small
\vspace*{-2em}
\infrule
  {
\jsubtype[\denv]{C_2}{C_3}  \andalso
    \jmatch{C_3}{\ob{\Bottom \extends P \extends \Any}}{P}{\bindcontra}{\biglistof{\listof{\Bottom}}}{\biglistof{\listof{C_3}}} \\[2pt]
    \jsolve{\ob{\Bottom \extends P \extends \Any}}{\biglistof{\listof{C_1}}}{\biglistof{\listof{C_3}}}{\subst{C_3}{P}}
  }
  { \jsolve{\ob{\Bottom \extends P \extends \Any, P \extends Q \extends \Any}}{\biglistof{\listof{C_1,C_2}}}{\biglistof{\listof{\Any,C_3}}}{\substxx{C_3}{P}{C_3}{Q}} }
\vspace*{-2em}
\end{figure}

\medskip
\noindent
As a result, the dispatch semipredicate on $d_2$ succeeds with
$\substxx{C_3}{P}{C_3}{Q}$, and the method call expression
is dispatched to the following instance of $d_2$:
\begin{center}
$D = \instantiate{d_2}{\substxx{C_3}{P}{C_3}{Q}}
= \mathit{append}(x \COLON \kwd{List}\ob{C_3}, y \COLON \kwd{List}\ob{C_3}) \COLON \kwd{List}\ob{C_3} = \dontcare$
\end{center}
It shows that the dynamic overloading resolution selects an
instance of the most specific applicable method declaration without any
ambiguity, and the return type of the chosen instance is a subtype of
the static return type, which means that the dynamic overloading
resolution preserves type.


\subsubsection{Correctness of Dispatch Semipredicate}\label{sec:prop}~\\
We prove the soundness and completeness of the dispatch semipredicate
by showing that each of its steps is sound and complete.

We first show that the type matching step is sound and complete.

The \textsf{match} rules compute sound bounds for method type parameters
by matching a ground type $g$ against an internal type $\alpha$ according
to a variance mark $V$.  That is, any substitution $\sigma$ satisfying
the bounds makes $\sigma\alpha$ a supertype, a subtype, or an equivalent type of
$g$ when $V$ is respectively $\bindco$, $\bindcontra$, or $\bindinv$.

\begin{lemma}(Soundness of Type Matching)\\
\label{lem:matchsound}
Let $\rep{K=\lobindup{\rep{\chi}}{P}{\rep{\eta}}}$,
$\jscope[\denv]{\obb{K}}$,
$\jwftype[\denv]{g}$,
$\jwftype[\denv\cup\setb{K}]{\alpha}$, and
$\sigma=\substb{g'}{P}$ with $\rep{\jwftype[\same\denv]{g'}}$.
% Suppose for all $\T$ and $\T'$,
% \begin{center}
% $\alpha\neq\T\context{\Bottom}$,
% $\rep{(\same\alpha=\T\context{\naryuniontypeb{\gamma}})\Rightarrow(\rep{\gamma\neq\T'\context{\same{P}}})}$, and
% $\rep{(\same\alpha=\T\context{\naryintersectiontypeb{\gamma}})\Rightarrow(\rep{\gamma\neq\T'\context{\same{P}}})}$.
% \end{center}

\begin{tabular}{l@{~}l}
If &
$\jmatch{g}{\obb{K}}{\alpha}{V}{\biglistofb{\listofb{\lbound}}}{\biglistofb{\listofb{\ubound}}}$,\\
&
$\jobeys[\denv]{\sigma}{\listofb{P}}{\listofb{\lbound_{\same{i}}}}{\listofb{\ubound_{\same{i}}}}$
for some $1\leq i\leq \countofb{\listofb{\lbound}}=\countofb{\listofb{\ubound}}$, and\\
&
$\jobeys[\denv]{\sigma}{\listofb{P}}{\biglistofb{\bigtypeunion \setb{\chi}}}{\biglistofb{\bigtypeintersection \setb{\eta}}}$,\\
then &
$\jcompare[\denv]{g}{\sigma\alpha}{V}$.
\end{tabular}
\end{lemma}

\begin{proof}
See Appendix~\ref{app:proofs}.
\end{proof}

Suppose an internal type $\alpha$ do not contain $\Bottom$ and any union
types and intersection types it contains do not contain any type
parameters. For a ground type $g$ and a variance mark $V$, if a
substitution $\sigma$ on $\listofb{P}$ makes $\sigma\alpha$ a supertype, a subtype, or an
equivalent type of $g$ when $V$ is respectively $\bindco$, $\bindcontra$, or
$\bindinv$, then the \textsf{match} rules find bounds of $\listofb{P}$
that $\sigma$ satisfies.

\begin{lemma}(Completeness of Type Matching)\\
Let $\rep{K=\lobindup{\rep{\chi}}{P}{\rep{\eta}}}$,
$\jscope[\denv]{\obb{K}}$,
$\jwftype[\denv]{g}$,
$\jwftype[\denv\cup\setb{K}]{\alpha}$, and
$\sigma=\substb{g'}{P}$ with $\rep{\jwftype[\same\denv]{g'}}$.
Suppose for all $\T$ and $\T'$,
\begin{center}
$\alpha\neq\T\context{\Bottom}$,
$\rep{(\same\alpha=\T\context{\naryuniontypeb{\gamma}})\Rightarrow(\rep{\gamma\neq\T'\context{\same{P}}})}$, and
$\rep{(\same\alpha=\T\context{\naryintersectiontypeb{\gamma}})\Rightarrow(\rep{\gamma\neq\T'\context{\same{P}}})}$.
\end{center}
\begin{tabular}{l@{~}l}
If &
$\jcompare[\denv]{g}{\sigma\alpha}{V}$
and
$\jobeys[\denv]{\sigma}{\listofb{P}}{\biglistofb{\bigtypeunion \setb{\chi}}}{\biglistofb{\bigtypeintersection \setb{\eta}}}$,\\
then &
$\jmatch{g}{\obb{K}}{\alpha}{V}{\biglistofb{\listofb{\lbound}}}{\biglistofb{\listofb{\ubound}}}$ and\\
&
$\jobeys[\denv]{\sigma}{\listofb{P}}{\listofb{\lbound_{\same{i}}}}{\listofb{\ubound_{\same{i}}}}$
for some $1\leq i\leq \countofb{\listofb{\lbound}}=\countofb{\listofb{\ubound}}$.
\end{tabular}
\label{lem:matchcomp}
\end{lemma}
\begin{proof}
See Appendix~\ref{app:proofs}.
\end{proof}

A subtle point in Lemma~\ref{lem:matchcomp} is the assumptions on type $\alpha$.
While the syntax of the \textsf{match} rules allows any internal types,
each \textsf{match} rule is used only with $\alpha$ coming from method
declarations.  The explicitly annotated types in method declarations are
external types, which do not contain $\Bottom$, union types, nor
intersection types.  On the contrary, inferred types for method type
parameters may contain such internal types.  For instance, in the
following example (in which type annotations appear for the two method calls):
\[
\begin{array}{lll}
\kwd{object}\;O'\ob{Q} & m'(x \COLON Q) \COLON Q = \dontcare & \kwd{end}\\
\kwd{object}\;O & m\ob{P}(x \COLON P, y \COLON P) \COLON P = O'\ob{P}.m'(x) \COLON P &
\kwd{end}\\
\lefteqn{O.m(\kwd{3}, \kwd{true}) \COLON \uniontype{\kwd{Int}}{\kwd{Boolean}}}
\end{array}
\]
the method type parameter $P$ of $m$ is inferred as
$\uniontype{\kwd{Int}}{\kwd{Boolean}}$ by the dispatch semipredicate
at run time.  Since this inferred type then flows into the method
declaration of $m'$, the ilk and the static type of method invocation
expression of $m'$ are matched against the union type.  However, because
our type system requires that both static and dynamic type inference
algorithms instantiate type parameters with non-$\Bottom$ types,
$\alpha$ does not contain $\Bottom$.  In addition, when $\alpha$
contains union types or intersection types, all the type parameters in
them are already instantiated with ground types.


Next, we show the soundness and completeness of the solving step.

The \textsf{solve} rules compute a valid substitution that satisfies
both the declared bounds and computed bounds of method type parameters.

\begin{lemma}(Soundness of Solving Bounds)\\
Let $\rep{K=\lobindup{\rep{\chi}}{P}{\rep{\eta}}}$ and
$\jscope[\denv]{\obb{K}}$.

\begin{tabular}{l@{~}l}
If &
$\jsolve{\obb{K}}{\biglistofb{\listofb{\lbound}}}{\biglistofb{\listofb{\ubound}}}{\sigma}$,\\
then &
$\jobeys[\denv]{\sigma}{\listofb{P}}{\listofb{\lbound_{\same{i}}}}{\listofb{\ubound_{\same{i}}}}$
for some $1\leq i\leq \countofb{\listofb{\lbound}}=\countofb{\listofb{\ubound}}$
and\\
&$\jobeys[\denv]{\sigma}{\listofb{P}}{\listofb{\bigtypeunion \setb{\chi}}}{\listofb{\bigtypeintersection \setb{\eta}}}$.
\end{tabular}
\label{lem:solvesound}
\end{lemma}
\begin{proof}
See Appendix~\ref{app:proofs}.
\end{proof}

If a valid substitution for method type parameters exists,
then the \textsf{solve} rules find one.

\begin{lemma}(Completeness of Solving Bounds)\\
Let $\rep{K=\lobindup{\rep{\chi}}{P}{\rep{\eta}}}$,
$\jscope[\denv]{\obb{K}}$, and
$\sigma=\substb{g}{P}$ with $\rep{\jwftype[\same\denv]{g}}$.

\begin{tabular}{l@{~}l}
If &
$\jobeys[\denv]{\sigma}{\listofb{P}}{\listofb{\bigtypeunion \setb{\chi}}}{\listofb{\bigtypeintersection \setb{\eta}}}$
and
$\jobeys[\denv]{\sigma}{\listofb{P}}{\listofb{\lbound}}{\listofb{\ubound}}$\\
& for some $\rep{\lbound}$ and $\rep{\ubound}$,\\
then &
$\jsolve{\obb{K}}{\biglistofb{\listofb{\lbound'}}}{\biglistofb{\listofb{\ubound'}}}{\sigma'}$ for some $\sigma'$\\
&
where $\listofb{\lbound}=\listofb{\lbound'_{\same{i}}}$ and $\listofb{\ubound}=\listofb{\ubound'_{\same{i}}}$
for some $1\leq i\leq \countofb{\listofb{\lbound'}}=\countofb{\listofb{\ubound'}}$.
\end{tabular}
\label{lem:solvecomp}
\end{lemma}
\begin{proof}
See Appendix~\ref{app:proofs}.
\end{proof}

Finally, we prove that the dispatch semipredicate is sound and complete.

If the dispatch semipredicate on a method declaration $d$ for an ilk
$k$ and a static return type $g$ succeeds,
then $d$ is dispatchable to $k$ and $g$.

\begin{theorem}(Soundness of Dispatch Semipredicate)\\
Let\ \
$\jwftype[\denv]{k}$,
$\jwftype[\denv]{g}$,
$\rep{K=\lobindup{\rep{\chi}}{P}{\rep{\eta}}}$,
$\jscope[\denv]{\obb{K}}$,
$\arrow(d)=\foralltypeb{K}{\arrowtype{\alpha}{\rho}}$,
$\jwftype[\denv\cup\setb{K}]{\alpha}$, and
$\jwftype[\denv\cup\setb{K}]{\rho}$.

\begin{tabular}{l@{~}l}
If &
$\jfancydispatch{k}{g}{d}{\sigma}$,\\
then &
$\jobeys[\denv]{\sigma}{\listofb{P}}{\listofb{\bigtypeunion \setb{\chi}}}{\listofb{\bigtypeintersection \setb{\eta}}}$,
$\jsubtype[\denv]{k}{\sigma\alpha}$,
and
$\jsubtype[\denv]{\sigma\rho}{g}$.
\end{tabular}
\label{thm:dispatchsound}
\end{theorem}
\begin{proof}
See Appendix~\ref{app:proofs}.
\end{proof}

If a method declaration $d$ is dispatchable to an ilk $k$ and a static
return type $g$, then the dispatch semipredicate on $d$ for $k$ and $g$
succeeds.

\begin{theorem}(Completeness of Dispatch Semipredicate)\\
Let $\rep{K=\lobindup{\rep{\chi}}{P}{\rep{\eta}}}$,
$\jscope[\denv]{\obb{K}}$,
$\sigma=\substb{g'}{P}$ with $\rep{\jwftype[\same\denv]{g'}}$,
$\jwftype[\denv]{k}$,
$\jwftype[\denv]{g}$,
$\arrow(d)=\foralltypeb{K}{\arrowtype{\alpha}{\rho}}$,
$\jwftype[\denv\cup\setb{K}]{\alpha}$, and
$\jwftype[\denv\cup\setb{K}]{\rho}$.

\begin{tabular}{l@{~}l}
If &
$\jobeys[\denv]{\sigma}{\listofb{P}}{\listofb{\bigtypeunion \setb{\chi}}}{\listofb{\bigtypeintersection \setb{\eta}}}$,
$\jsubtype[\denv]{k}{\sigma\alpha}$,
and
$\jsubtype[\denv]{\sigma\rho}{g}$,\\
then &
$\jfancydispatch{k}{g}{d}{\sigma'}$ for some $\sigma'$.
\end{tabular}
\label{thm:dispatchcomp}
\end{theorem}
\begin{proof}
See Appendix~\ref{app:proofs}.
\end{proof}


\subsection{Type Soundness of FGFV}\label{sec:sound}
We now prove the type soundness of FGFV.
First, we show that for any applicable method declaration instance
for a given method invocation, we can find an instance of
the most specific method declaration whose return type is
a subtype of the return type of the applicable instance.
\begin{lemma}(Existence of a Type-safe Most Specific Method Instance)\\
\label{lem:instance}
Let $\setb{d_0'}=\visible{\tuple{C,m}}$ for $\jwftype{C}$, and
$\tuple{d_0,\setb{\sigma_0}}\in\applicable{\setb{d_0'}}{\alpha}$
and
$\tuple{d,\setb{\sigma}}\in\mostspecific{\applicable{\setb{d_0'}}{\alpha}}$
for $\jwftype{\alpha}$ and $\alpha\neq\Bottom$.

\begin{tabular}{l@{~}l}
If & $\,D_0\in \bigsetb{\instantiate{\same{d_0}}{\sigma_0}}$
with
$\arrow(D_0)=\foralltype{}{\arrowtype{\dontcare}{\rho_0}}$,\\
then &
$\,\exists\;D\in\bigsetb{\instantiate{\same{d}}{\sigma}}$
with
$\arrow(D)=\foralltype{}{\arrowtype{\dontcare}{\rho}}$
such that $\jsubtype{\rho}{\rho_0}$.
\end{tabular}
\end{lemma}
\begin{proof}
See Appendix~\ref{app:proofs}.
\end{proof}

Next, we show that no matter which applicable method instance we choose
for a method invocation expression at compile time, there exists
a type-safe most specific method instance at run time.
\begin{lemma}(Existence of a Type-safe Most Specific Method Instance at Run Time)\\
\label{lem:runtimeinstance}
Let
$\setb{d_0'}=\visible{\tuple{C,m}}$ and $\setb{d'}=\visible[\denv]{\tuple{G,m}}$
for $\jwftype{C}$, \mbox{$\jwftype[\denv]{G}$}, and $\Delta = \denv \cup \set{\rep{\lobindup{\rep{\chi}}{P}{\rep{\eta}}}}$.
Also, let
$\tuple{d_0,\setb{\sigma_0}}\in\applicable{\setb{d_0'}}{\alpha}$
for \mbox{$\jwftype{\alpha}$} and $\alpha\neq\Bottom$, and
$\tuple{d,\setb{\sigma}}\in\mostspecific[\denv]{\applicable[\denv]{\setb{d'}}{k}}$
for $\jwftype[\denv]{k}$
with
\mbox{$\jsubtype[\denv]{G}{\sigma C}$} and $\jsubtype[\denv]{k}{\sigma\alpha}$
for
$\sigma=\substb{g'}{P}$ with
$\jobeys[\denv]{\sigma}{\listofb{P}}{\listofb{\bigtypeunion \bdb{\chi}}}{\listofb{\bigtypeintersection \bdb{\eta}}}$
and $\rep{\jwftype[\same{\denv}]{g'}}$.

\begin{tabular}{l@{~}l}
If &
$D_0\in\bigsetb{\instantiate{\same{d_0}}{\sigma_0}}$
with
$\arrow(D_0)=\foralltype{}{\arrowtype{\dontcare}{\rho_0}}$
\\
then &
$\exists\;D\in\bigsetb{\instantiate{\same{d}}{\sigma}}$
with
$\arrow(D)=\foralltype{}{\arrowtype{\dontcare}{h}}$
such that $\jsubtype[\denv]{h}{\sigma\rho_0}$.
\end{tabular}
\end{lemma}
\begin{proof}
See Appendix~\ref{app:proofs}.
\end{proof}

We now show that if a type-safe most specific method instance exists,
the \textit{firstDispatchable} rule finds one. This lemma follows from
the soundness and completeness of the dispatch semipredicate.

\begin{lemma}(Validity of the firstDispatchable Rule)\\
\label{lem:fstdisp}
Let $\setb{d'}=\visible[\denv]{\tuple{G,m}}$ for $\jwftype[\denv]{G}$.

\begin{tabular}{l@{~}l}
If &
$\tuple{d,\setb{\sigma}}\in\mostspecific[\denv]{\applicable[\denv]{\setb{d'}}{k}}$
for $\jwftype[\denv]{k}$
and
$D\in\bigsetb{\instantiate{\same{d}}{\sigma}}$\\
& with
$\arrow(D)=\foralltype{}{\arrowtype{\dontcare}{h}}$
and
$\jsubtype[\denv]{h}{g}$ for $\jwftype[\denv]{g}$,
\\
then &
$\firstdispatchable[\denv]{\topologicalsort[\denv]{\setb{d'}}}{k}{g}=D_1$\\
& with
$D_1\in\bigsetb{\instantiate{\same{d}}{\sigma}}$,
$\arrow(D_1)=\foralltype{}{\arrowtype{h_1}{h_2}}$,
$\jsubtype[\denv]{k}{h_1}$,
and
$\jsubtype[\denv]{h_2}{g}$.
\end{tabular}
\end{lemma}
\begin{proof}
See Appendix~\ref{app:proofs}.
\end{proof}

Finally, we prove the type soundness of FGFV.
Since the main focus of this paper is resolution of overloaded method
invocation and the other cases are conventional, we prove only the
method invocation case: a well-typed method invocation
can be reduced via \rncite{R-Method} and the reduced expression preserves
the type of the method invocation.

\begin{theorem}(Type Soundness of Method Invocation Expression)\\
\label{thm:soundinvok}
Let
$\listofb{d} = \topologicalsort[\denv]{\visible[\denv]{\tuple{O\obb{g''},m}}}$
for 
\mbox{$\objectdecl{O}{\dontcares}{\rep{x\COLON\dontcare}}{\dontcares}{\dontcares}\in\denv$} and
$\jwftype[\denv]{O\obb{g''}}$. 

\begin{tabular}{l@{~}l}
If &
$\,\jtype[\denv;\emptyseq]{O\obb{g''}(\bar{v}).m(\rep{v'})}{\tuple{\dontcare,g}}$,\\
then &
$(\dontcare, \methoddef{m}{}{\rep{x'\COLON \dontcare}}{\dontcare}{\epsilon})=\firstdispatchable[\denv]{\listofb{d}}{\tuple{O\obb{g''},\bar{\ilk{v'}}}}{g}$\\
and &
\mbox{$\jtype[\denv;\emptyseq]{\subst{O\obb{g''}(\bar{v})}{\kwd{self}}\substb{v}{x}\substb{v'}{x'}\epsilon}{\tuple{\dontcare,g'}}$}
with
$\jsubtype[\denv]{g'}{g}$.
\end{tabular}
\end{theorem}
\begin{proof}
See Appendix~\ref{app:proofs}.
\end{proof}

\begin{theorem}(Type Soundness of FGFV)\\
\label{thm:typesoundness}
If
$\,\jtype[\denv;\emptyseq]{e}{\tuple{\epsilon,g}}$
and
$\jevalstar[\denv]{\epsilon}{v}$, then
$\jtype[\denv;\emptyseq]{v}{\tuple{\dontcare,g'}}$
and
$\jsubtype[\denv]{g'}{g}$.
\end{theorem}
\begin{proof}
See Appendix~\ref{app:proofs}.
\end{proof}

\subsection{Discussion}\label{sec:discuss}
Defining a provable system for an object-oriented language requires addressing
various technical challenges due to the structural asymmetries inherent
in the language design.  One can have a variable of type $\Any$,
but one cannot have a variable of type $\Bottom$.  The ilk of an actual
argument must be a subtype, not a supertype, of the method's parameter
type.  The ilk of a returned value must be a subtype, not a supertype, of
the method's return type.  A trait may declare ancestors but not
descendants.  Such asymmetries prohibit
seemingly symmetric features from being treated as duals.  One might think
at first glance that upper bounds and lower bounds of type parameters
should be related symmetrically (perhaps treated as duals), but they cannot be.

For example, the ``Ancestor Meet Rule''
(expressed in this paper by rule \rncite{Anc-Same-Trait}) was introduced
into Fortress in 2012 by Guy Steele and David Chase~\cite{dchase12} to
design an expressive type system with union and intersection types.  This rule
permits proofs that (a)~if $T\ob{P}$ is covariant in $P$ then
$T\ob{\intersectiontype{\alpha}{\alpha'}}\equiv \intersectiontype{T\ob{\alpha}}{T\ob{\alpha'}}$, and
(b)~if $T\ob{P}$ is contravariant in $P$ then
$T\ob{\uniontype{\alpha}{\alpha'}}\equiv \intersectiontype{T\ob{\alpha}}{T\ob{\alpha'}}$
(that last set operator is indeed $\sqcap$, not $\sqcup$).
These equivalences allow simplification of type expressions by moving type union
and type intersection operators out of type arguments.
However, the ``obvious other two cases'' are not available.
When $T\ob{P}$ is covariant in $P$, $T\ob{\uniontype{\alpha}{\alpha'}}$
is not equivalent to some simple combination of $T\ob{\alpha}$ and $T\ob{\alpha'}$;
the best one can say is that $\uniontype{T\ob{\alpha}}{T\ob{\alpha'}} \subtypeof T\ob{\uniontype{\alpha}{\alpha'}}$.

In 2012, the developers of Fortress explored an alternate set of language restrictions
that exploited these asymmetries to ensure that exponential blowup could occur at run time
only when union types interacted with contravariance (the dual case of intersection types
interacting with covariance could not occur).  They then
contemplated two alternatives for avoiding the exponential blowup.
One possibility is to eliminate contravariant type parameters and contravariant type relationships.
This has the consequence that arrow types must be treated as invariant, rather than contravariant,
in the type to the left of the arrow.
This could make it inconvenient to use higher-order functions such as
$\mathit{map}$, but one can compensate by introducing rules of
coercion, so that if a function of type $\arrowtype{\alpha}{\rho}$ is
used in a situation where a function of type $\arrowtype{\alpha'}{\rho'}$ is required,
it can be coerced to a new function object if $\alpha'\subtypeof\alpha$ and $\rho\subtypeof\rho'$.
The other possibility is to eliminate union types from the type
theory; this might require restricting the language so that, for
example, in an \texttt{if}-\texttt{then}-\texttt{else}
expression, the type of the $\texttt{then}$ subexpression must be
either a supertype or a subtype of the $\texttt{else}$ subexpression.
Alternatively, some sort of pseudo-join (least common existing supertype) operation could be used.
One could go even further, also eliminating intersection types from the type theory
by further restricting the language and possibly using a pseudo-meet operation.

% In this paper, we carefully designed the dynamic dispatch algorithm to
% make a balanced tradeoff between static expressivity and dynamic
% efficiency.  We accept the potential for exponential blowup during
% dynamic dispatch because we believe that such cases are unlikely to
% arise in practice and the use of union and intersection types in the
% type theory, even if not directly expressible by the programmer, makes
% the language more expressive.  Also, while one could in principle
% allow any lower bound to refer to any type parameter, dynamic
% overloading resolution would then have to solve a general set of type
% inequalities, requiring a complex technique with multiple iterations;
% we chose instead to impose left-to-right scoping on method type
% parameters, so as to allow the use of a simple, one-pass right-to-left scan.

We carefully designed our dynamic dispatch algorithm to
make a balanced tradeoff between static expressivity and dynamic
efficiency.  We accept the potential for exponential blowup
% during dynamic dispatch
because we believe that such cases are unlikely to
arise in practice and the use of union and intersection types in the
type theory, even if not directly expressible by the programmer, makes
the language more expressive.  Also, while one could in principle
allow any lower bound to refer to any type parameter, dynamic
overloading resolution would then have to solve a general set of type
inequalities, requiring a complex technique with multiple iterations;
we chose instead to impose left-to-right scoping on method type
parameters, allowing use of a simple, one-pass right-to-left scan.
We believe this tradeoff represents a relatively small sacrifice of expressivity for
a large gain in dynamic efficiency.

Dynamically inferring type instantiations can incur run-time overhead.
In simple invariant cases, we expect the run-time cost to be relatively low but nonzero.
For example, if we pass a value of type \texttt{List$\ob{\texttt{String}}$} to a method such as
\textit{foo}$\ob{T}$\texttt{($x$:List$\ob{T}$)},
$T$ is matched only once and no union or intersection types are involved.
If \texttt{List} is invariant in $T$, it will come down to extracting the type \texttt{String}
from the type \texttt{List$\ob{\texttt{String}}$} and binding $T$ to that.
For more costly covariant and contravariant cases,
as in certain JVM implementations, such techniques as caching or dynamic recompilation
in reaction to the relative frequencies of types actually encountered may help.

The biggest challenge we faced in trying to find a provable system was
grappling with the language asymmetries described above and trying to
understand their consequences.  In particular, the Fortress team had
assumed, because it seemed from several working examples to be in accord
with intuition, that once the tightest possible upper and lower bounds
for a method type parameter have been found, then the union
(or pseudo-union) of the lower bounds should be chosen as the type to be used for that type parameter.
% But we were unable to complete the proofs of correctness presented here until we abandoned that assumption.
% We then considered the hybrid technique of using the intersection of upper bounds only when
% the union of lower bounds is \Bottom, but we were unable to get that to work either.
% That is why the dynamic dispatch algorithm presented here always chooses the intersection of the upper bounds.
% It remains an open question whether there is a provably correct dynamic dispatch technique that
% chooses the union of the lower bounds.
We considered this and the hybrid technique of using the intersection of upper bounds only when
the union of lower bounds is \Bottom, but we were unable to
complete the proofs of correctness of both the techniques.
Since the minimum requirement for a type to instantiate a method type parameter
is that it satisfies all the lower and upper bounds of the type parameter,
the dynamic dispatch algorithm presented here always chooses
the intersection of the upper bounds.

In the work presented here, as in previous work, one of the most difficult parts
is finding the necessary substitution $\sigma$ to satisfy the rule \rncite{E-Sub} for existential
subtyping during static analysis.  This is more difficult than the corresponding calculation
during dynamic dispatch because at run time one always matches an ilk (which contains no
type parameters) to a type, but at compile time one must match a type against a type,
where each type may contain type parameters.
Indeed, we do not present a specific static type inference
algorithm that finds the substitution $\sigma$ in rule \rncite{E-Sub},
but only refer to previous work in this area.
However, we conjecture that the \textsf{match} rule presented here can
serve this role after being extended to handle the case of matching
two type parameters.
